% Generated draft: requires mathematical and source review.
\documentclass[11pt]{article}
\usepackage[margin=25mm]{geometry}
\usepackage{fontspec}
\setmainfont{Latin Modern Roman}
\usepackage{amsmath,amssymb,mathtools,unicode-math}
\setmathfont{Latin Modern Math}
\usepackage{enumitem,needspace,xurl,hyperref}
\setlength{\parindent}{0pt}
\setlength{\parskip}{5pt}
\setlength{\emergencystretch}{3em}
\newcommand{\N}{\mathbb{N}}
\newcommand{\Z}{\mathbb{Z}}
\newcommand{\Q}{\mathbb{Q}}
\newcommand{\R}{\mathbb{R}}
\newcommand{\C}{\mathbb{C}}
\newcommand{\F}{\mathbb{F}}
\newcommand{\A}{\mathbb{A}}
\newcommand{\PP}{\mathbb{P}}
\newcommand{\E}{\mathbb{E}}
\newcommand{\eps}{\varepsilon}
\newcommand{\dd}{\,\mathrm{d}}
\begin{document}
\section*{UnsolvedMath ID: 7200087}
\textbf{Source problem code:} catalogue entry
\subsection{Definitions and mathematical statement}
In the usual closed polyhedral surface interpretation, seek a nonconvex polyhedron \(P\subset\R^3\) with an embedded boundary surface, so there are no self-intersections. Let \(\mathcal F(P)\) be its finite collection of polygonal faces. The requested conditions are
\[
 |\mathcal F(P)|>7,\qquad
 \forall F_1\ne F_2\in\mathcal F(P),\quad
 F_1\cap F_2\text{ contains a common edge of }P.
\]
Equivalently, the graph with one vertex for each face and one adjacency for each edge-sharing pair is complete. Sharing only a vertex is insufficient. The condition applies to all faces of the polyhedron, not only to a selected subcollection.

The source specifies no topology, so a spherical surface must not be imposed: higher genus is allowed unless the missing convention excludes it. It also does not specify whether all faces are simple planar disks, whether any two faces may share more than one edge, or which polyhedral manifold conventions are required. These choices need clarification before an exact existence target can be fixed.

The research summary's assertion of a convex maximum of seven faces and its identification with the Cs\'asz\'ar polytope are not reliable formulations of the face question. They may conflate vertices with faces and convex with nonconvex examples. In particular the face adjacency graph of a convex 3-polytope is planar, so a complete face adjacency claim with seven vertices is incompatible with that standard setting. No maximum or uniqueness assertion from that summary is adopted.
\subsection{Short English statement}
Does an embedded nonconvex polyhedron exist with at least eight faces such that every pair of distinct faces shares an edge? The topology is unrestricted in the source, and the summary's convex and Cs\'asz\'ar claims need correction.
\subsection{Sources}
\begin{itemize}
\item Local UnsolvedMath registry, permanent ID 7200087, source code catalogue entry; updated 2026-08-06; no proposer or proposal year supplied.
\item Wikipedia, \emph{List of unsolved problems in mathematics}, Geometry bullet 87, pinned revision 1366636767: \url{https://en.wikipedia.org/w/index.php?title=List_of_unsolved_problems_in_mathematics&oldid=1366636767#Geometry}. Registry access date 2026-07-29; MediaWiki extraction. Wikipedia text is CC BY-SA 4.0, with attribution and ShareAlike handling required for release. The provenance states unpublished pending primary-source review, despite the registry published flag.
\item The summary mentions Cs\'asz\'ar and Gimbel without precise references, and expressly marks the nonconvex question's status uncertain. No verified primary citation is supplied.
\end{itemize}
\end{document}
